\section{Theoretical Setup}\label{sec:setup}

All logarithms denoted by \(\log\) are natural.  For \(t\ge1\), define
\[
\log^*t
=\min\{r\in\mathbb N_0:\log_2^{(r)}(t)\le1\},
\]
where \(\log_2^{(0)}(t)=t\) and \(\log_2^{(r)}\) is the \(r\)-fold
iterate of the base-two logarithm for \(r\ge1\).

Let \((X,\Sigma)\) be a measurable instance space and
\(C\subseteq\{0,1\}^X\) a binary concept class.  A partition
\(\mathcal P\) of \(X\) is \emph{Cartesian for \(C\)} if, after writing
\(C_B=\{c|_B:c\in C\}\), the restriction map
\[
c\longmapsto(c|_B)_{B\in\mathcal P}
\]
is a bijection from \(C\) onto \(\prod_{B\in\mathcal P}C_B\).  A
Cartesian partition is finest if it refines every Cartesian partition;
such a partition is canonical up to relabeling.

\begin{assumption}[Canonical disjoint Cartesian factorization]
\label{assump:canonical-product}
The class \(C\) admits a finite finest Cartesian partition
\[
X=\bigsqcup_{i=1}^kX_i,\qquad 1\le k<\infty,
\]
of the whole domain, and the restriction map identifies \(C\) with the
full product
\[
C=\prod_{i=1}^k C_i,
\qquad
C_i=\{c|_{X_i}:c\in C\}.
\]
This is a static property of \(C\), independent of any learner,
dataset, distribution, event, or generated object.
\end{assumption}

\begin{assumption}[Nonconstant finite-Littlestone VC-one factors]
\label{assump:vc-one-factors}
Every \(C_i\) is nonconstant, has
\(\operatorname{VC}(C_i)=1\), and has finite
\[
d_i=\operatorname{LD}(C_i).
\]
No ordering, finite-domain, finite-cardinality, properness, or
computational assumption is imposed.
\end{assumption}

Under Assumptions~\ref{assump:canonical-product} and
\ref{assump:vc-one-factors}, define
\begin{equation}\label{eq:setup-complexities}
s_i=1+\log^*(d_i+1),
\qquad
M_\oplus(C)=\sum_{i=1}^k s_i,
\qquad
\omega_i=\frac{s_i}{M_\oplus(C)}.
\end{equation}
For a candidate sample size \(n\in\mathbb N=\{1,2,\ldots\}\), define
the factor budget
\begin{equation}\label{eq:setup-factor-budget}
m_{n,i}=\max\{8,\lceil4n\omega_i\rceil\}.
\end{equation}

For \(x,x'\in X_i\), write
\[
x\equiv_i x'
\quad\Longleftrightarrow\quad
c_i(x)=c_i(x')\ \text{for every }c_i\in C_i.
\]
Let \(Q_i=X_i/{\equiv_i}\) and let
\(\kappa_i:X_i\to Q_i\) be the quotient map.  Give \(Q_i\) its discrete
sigma-field \(2^{Q_i}\), and let
\[
\Sigma_i=\{A\cap X_i:A\in\Sigma\}.
\]

\begin{assumption}[Finite-or-countable measurable evaluation quotient]
\label{assump:countably-coded-evaluation}
Every \(X_i\) belongs to \(\Sigma\), every evaluation quotient \(Q_i\)
is finite or countable, and every quotient cell
\(\kappa_i^{-1}(\{q\})\), \(q\in Q_i\), belongs to \(\Sigma_i\).
Equivalently,
\(\kappa_i:(X_i,\Sigma_i)\to(Q_i,2^{Q_i})\) is measurable.
This is a primitive static condition on the evaluation structure.  It
does not assume any data-dependent version space, core, depth, selector,
learner, or output kernel to exist or be measurable, and it does not
require \(X_i\) to be finite, countable, countably generated, or
standard Borel.
\end{assumption}

Each \(c_i\in C_i\) has a unique representative
\(\bar c_i:Q_i\to\{0,1\}\) satisfying
\(c_i=\bar c_i\circ\kappa_i\).  Put
\[
\bar C_i=\{\bar c_i:c_i\in C_i\},
\qquad
\mathcal H_i=\{0,1\}^{Q_i},
\]
where \(\mathcal H_i\) has the product Borel sigma-field
\(\mathscr H_i\) generated by finite evaluation cylinders.  Also put
\[
\mathcal H^\oplus=\prod_{i=1}^k\mathcal H_i,
\qquad
\mathscr H^\oplus=\bigotimes_{i=1}^k\mathscr H_i.
\]
For \(\bar h=(\bar h_1,\ldots,\bar h_k)\in\mathcal H^\oplus\), define
the decoded hypothesis
\begin{equation}\label{eq:setup-decoder}
h_{\bar h}(x)=\bar h_i(\kappa_i(x))
\qquad(x\in X_i).
\end{equation}
Assumption~\ref{assump:countably-coded-evaluation} makes every target
and every decoded hypothesis \(\Sigma\)-measurable.

\begin{assumption}[Approximate-DP parameter range]
\label{assump:global-privacy-range}
The privacy parameters satisfy
\[
0<\varepsilon\le\frac1{10},
\qquad
0<\delta<1.
\]
The upper theorem has no further restriction on \(\delta\).
\end{assumption}

The lower argument uses the following primitive numerical condition
only at the particular candidate to which it is applied.

\begin{assumption}[Candidate-wise lower-bound budget]
\label{assump:candidate-delta-budget}
For the candidate \(n\in\mathbb N\),
\begin{equation}\label{eq:setup-candidate-delta-budget}
0<\delta\le
\min\left\{
\frac1{n\log(n+1)},
\min_{1\le i\le k}
\frac{c_\delta}
{m_{n,i}^2\log(m_{n,i}+1)}
\right\},
\end{equation}
where \(c_\delta>0\) is the universal constant in the unrestricted
one-factor lower interface of
\citet{alonLivniMalliarisMoran2019}.  This condition is checked
separately at the actual candidate; it is not a uniform-in-\(n\)
schedule or a condition on a realized object.
\end{assumption}

Our objective is a conditional two-sided direct-sum theorem for this
measurable Cartesian subclass.  The upper result will hold throughout
Assumption~\ref{assump:global-privacy-range}; the lower result will be
candidate-wise under Assumption~\ref{assump:candidate-delta-budget}.
The theorem therefore does not purport to characterize arbitrary
finite-Littlestone classes or classes with uncountably many evaluation
types.
